% Français 9115, batch 2 — Gallica views 21–40.
% Pass: Opus 5 (claude-opus-5), 2026-09-12 — first pass, unchecked against
% the views by a human.
%
% What the twenty views hold. Batch 2 is the direct continuation of batch 1,
% and like it holds one piece and nothing else: the fair copy of « Exposition
% des principes qui servent de base à la théorie des surfaces élastiques ».
% Twenty consecutive written leaves, no blanks, no guard leaves, the same
% regular hand throughout, the text running from view to view without a break.
% It enters at the middle of the sentence view 20 broke off in, and it leaves
% off in the middle of another at the foot of view 40, which is again the
% batch boundary and not the end of the piece. Batch 3 continues it.
%
% What changes between the two batches is the register. Batch 1 was prose —
% the Avertissement, the two genres of hypothesis, the a priori examination
% of the two candidates, argued in words with the occasional term. Here the
% argument turns analytic: the variational calculation is carried out for
% each hypothesis in turn, the integrations by parts are done, the boundary
% terms are separated from the interior ones, and the equation of the
% vibrating plate is reached and written down. The prose that survives is
% about what the calculation is for, and it is still the best of it.
%
%   v. 21–24   End of N° 6 and the whole of N° 7. The comparison of the two
%              hypotheses, finished: the first is general like the definition
%              of the elastic forces themselves and applies to every natural
%              figure the surface can take; the second can only apply to the
%              single case where 1/ρ + 1/ρ′ − (1/R + 1/R′) and 1/ρ − 1/ρ′
%              vanish together, which is to say to naturally plane surfaces.
%              N° 7 then establishes what separates the two analytically:
%              the terms they give differ by the exact variation
%              δ∬ k dx dy/ρρ′, which is nothing for a plane surface with
%              fixed extremities, and is not nothing otherwise.
%   v. 25–27   The long quotation from the « savant auteur du mémoire sur les
%              surfaces élastiques », N° 14 p. 202 of his memoir, on the
%              repulsive forces and the limits of the integrals; then her
%              answer. The most personal page of the batch is view 27: « Je
%              ne me dissimule pas que le point de vue sous lequel j'ai
%              envisagé la question est entièrement neuf ; et c'est pour moi
%              un grand sujet de défiance », and it closes « je supplie les
%              géomètres de lui accorder quelqu'attention ».
%   v. 28      N° 8, on the maximum-or-minimum property of the integral, with
%              a second citation (N° 23 p. 221), and the announcement that
%              from here on only vibrating plates are treated.
%   v. 29–31   « Conclusions analytiques déduites de l'une et l'autre
%              hypothèses », N° 9 to N° 11. N° 9 takes her own term for a
%              naturally plane surface, reduces it to
%              −(r+t) δ(r+t) with r = d²z/dx², t = d²z/dy², and integrates by
%              parts. N° 10 does the same for the rival's term, by way of the
%              radical √{(r+t)² − 4(rt − s²)}. N° 11 collects what stays
%              under the double sign — the interior points — finds it the
%              same under both hypotheses, multiplies by the constant n² and
%              adds the term of the accelerating forces, and view 31 writes
%              the equation of the vibrating plate:
%              d²z/dt² + n²(d⁴z/dx⁴ + 2 d⁴z/dx²dy² + d⁴z/dy⁴) = 0,
%              with the remark, in her own words, that this is the equation
%              calculated on her hypothesis and found « dans la pièce que
%              j'ai envoyée pour le concours de janvier 1814 ».
%   v. 31–40   « Comparaison entre les résultats du calcul et les phénomènes
%              observés », N° 12 to N° 15. N° 12 the boundary terms and the
%              parallelogram plate; a first N° 13 (v. 32–34) on her own
%              hypothesis, the nodal lines and the free lines, and the
%              property that a plate cut along its nodal lines gives the same
%              sound in each piece; a second N° 13 (v. 35–37) on the rival's,
%              where the extra term 4∫ d²z/dxdy δ(dz/dy) dy does not vanish
%              and the two hypotheses part company; N° 14 the linear integral
%              and the invariable points; N° 15 (v. 38–40) the loaded plate —
%              wax applied to one edge — and the closing pages on the weight
%              of two celebrated names, breaking off at « que les géomètres
%              justement avares de leur tems ».
%
% The numbering on the leaves, and why no \folio{} is written. The two ink
% series batch 1 recorded continue through this batch without a break:
%
%   odd views  21 23 25 27 29 31 33 35 37 39  →  11 12 13 14 15 16 17 18 19 20
%   struck, same views                        →  19 21 23 25 27 29 31 33 35 37
%   even views 22 24 26 28 30 32 34 36 38 40  →  20 22 24 26 28 30 32 34 36 38
%
% Every one of those thirty numbers was read on the view; none is inferred.
% They fit batch 1's reconstruction exactly — leaf N on the recto, the struck
% number being that recto's cancelled odd page 2N−3, the verso carrying the
% page 2N−2 — but that reconstruction remains this project's and not the
% leaf's, the numbers are in ink and not in the Bibliothèque's pencil, and no
% pencilled foliation was found anywhere in the batch. So, as in batch 1, no
% \folio{} is written: the numbers stand verbatim in a \note{} at the head.
%
% Notation. Continues batch 1: ρ, ρ′ the principal radii of the strained
% figure, R, R′ those of the natural figure; δ the variation. The second and
% higher derivatives are written throughout with a straight d — d²z/dx²,
% d²z/dxdy, d⁴z/dx⁴, d⁴z/dx²dy² — where the modern notation is partial, and
% they are set here as she sets them. Two constants enter: k, in the
% integrals of N°s 7 and 8, and n², « proportionnel à l'épaisseur et a
% l'élasticité naturelle de la surface », in N° 11. ω and A carry the
% wavelength in N°s 13–14; m and n are the integers of the nodal figure; A
% and B the two constants of the linear integral. Her « ∫ » and her « δ » are
% drawn almost alike, a tall looped s; the reading here follows the sense,
% and the few places where the stroke alone would not decide are noted.
%
% What could not be read. Five heavy ink blots, each obliterating a single
% sign inside a bracket, are marked \ill{}: four of them (v. 33 twice, v. 35
% once, v. 30 once) stand between d²z/dx² and d²z/dy² where the sense wants a
% minus, and one (v. 37) stands between z and 0 where the sense wants an
% equals. One token on view 30 is struck and overwritten past reading. None
% of these was guessed. Her own underlining is set \emph{}; \uncertain{}
% renders as an underline too, so a reader of the PDF should take this
% header's word for which is which.
%
% Nothing here was checked against any published transcription.
\documentclass[11pt,a4paper]{article}
% The preamble shared by every transcription of the mathematicians' manuscripts in Gallica.
%
% It rests on one idea: **uncertainty compounds, it does not resolve.** A
% transcription that smooths over a doubtful word has destroyed the only thing
% it was made for — a reader must be able to tell, at every point, what was on
% the page and what was supplied. The apparatus macros below are therefore the
% critical apparatus, not decoration.
%
% The same commands are recognised by `scripts/render.mjs`, which produces the
% site's reading view, and by `scripts/tei.mjs`. A macro added here must be
% added there too, or rendering fails loudly — which is the wanted behaviour.
%
% Comments here are English, like the rest of the repository. The documents
% this preamble sets are French, and that is deliberate: see the skills.

\ProvidesPackage{germain}[2026/09/15 Transcriptions of mathematicians' manuscripts in Gallica]

\usepackage{amsmath,amssymb,amsthm}
\usepackage{mathrsfs}
% Kept from the parent preamble so the renderer's diagram support stays
% available; these manuscripts draw no commutative diagrams, and nothing here uses it.
\usepackage{tikz-cd}
\usepackage[english,french]{babel}
\usepackage{fontspec}

% --- Greek ------------------------------------------------------------------
% The text face stays fontspec's default, Latin Modern, which has no Greek:
% XeTeX drops every glyph it lacks with a line in the log and nothing on the
% page, and the Greek words the manuscripts quote vanished from the PDFs. So
% the Greek and Greek Extended blocks get a XeTeX character class of their own,
% and entering or leaving a run of it switches the family to CMU Serif — the
% same Computer Modern design, with polytonic Greek — and back. Only the
% family changes: a Greek word in \textit or \textbf keeps its shape and series.
% The transitions to and from every other class carry the tokens that class
% has towards an ordinary letter, so French spacing before « ; » or inside
% « » still holds after a Greek word. No grouping, so no brace or \endgroup
% can come out unbalanced; maths is left alone, since XeTeX inserts nothing
% there. Kept identical in `exercices.sty`.
\makeatletter
\newfontfamily\germainGreekfont{cmun}[
  Extension=.otf, NFSSFamily=germaingreek,
  UprightFont=*rm, ItalicFont=*ti, SlantedFont=*sl,
  BoldFont=*bx, BoldItalicFont=*bi, BoldSlantedFont=*bl]
\newXeTeXintercharclass\germain@greek
\def\germain@greekrange#1#2{%
  \@tempcnta=#1\relax
  \loop
    \XeTeXcharclass\@tempcnta=\germain@greek
    \ifnum\@tempcnta<#2\advance\@tempcnta\@ne\repeat}
\germain@greekrange{"0370}{"03FF}
\germain@greekrange{"1F00}{"1FFF}
\def\germain@greekprev{\rmdefault}
\def\germain@greekin{%
  \edef\germain@greekprev{\f@family}\fontfamily{germaingreek}\selectfont}
\def\germain@greekout{\fontfamily{\germain@greekprev}\selectfont}
\def\germain@greekjoin{%
  \XeTeXinterchartokenstate=\@ne
  \@tempcnta=\z@
  \loop
    \ifnum\@tempcnta=\germain@greek\else
      \XeTeXinterchartoks\germain@greek\@tempcnta=\expandafter{%
        \expandafter\germain@greekout\the\XeTeXinterchartoks\z@\@tempcnta}%
      \XeTeXinterchartoks\@tempcnta\germain@greek=\expandafter{%
        \the\XeTeXinterchartoks\@tempcnta\z@\germain@greekin}%
    \fi
    \ifnum\@tempcnta<4095\advance\@tempcnta\@ne\repeat}
% babel-french sets its own classes, and saves and restores the interchar
% state, each time a language is selected: join again after it does.
\addto\extrasfrench{\germain@greekjoin}
\addto\extrasenglish{\germain@greekjoin}
\AtBeginDocument{\germain@greekjoin}
\makeatother

\usepackage{geometry}
\geometry{a4paper,margin=25mm,marginparwidth=28mm}
\usepackage{marginnote}
\usepackage{xcolor}
\usepackage[normalem]{ulem}
\setlength{\emergencystretch}{3em}
\usepackage{hyperref}
\hypersetup{colorlinks=true,linkcolor=black,urlcolor=black,pdfborder={0 0 0}}

% --- Batch metadata ---------------------------------------------------------
% Taken as they stand by the rendering script, which shows them at the head of
% the reading view. They fix what the file claims to cover. The macro names are
% the parent project's, so that the scripts stay shared; \folder here means the
% volume's slug — `fr-9115` — and \pages counts Gallica views.

\newcommand{\folder}[1]{\gdef\grothFolder{#1}}
\newcommand{\batch}[1]{\gdef\grothBatch{#1}}
\newcommand{\pages}[2]{\gdef\grothFirst{#1}\gdef\grothLast{#2}}
\newcommand{\foldertitle}[1]{\gdef\grothFoldertitle{#1}}

% The holder's own shelfmark, printed as they write it: « Français 9115 ».
\newcommand{\shelfmark}[1]{\gdef\grothShelfmark{#1}}

% The dating, copied verbatim from the holder's catalogue — never inferred
% here. « XIXe siècle » is the BnF's; a pass that dates a leaf from its content
% says so in a \note{}, as its own inference.
\newcommand{\dating}[1]{\gdef\grothDating{#1}}

% The status notice, printed in the title block and repeated as a diagonal
% watermark on every page of the PDF. The manuscripts are in the public
% domain; what this line declares is not a right but a quality — an
% unverified first pass by a machine. The skills write the line; a file
% without it gets no watermark, so leaving it out is visible in the output.
\newcommand{\watermark}[1]{\gdef\grothWatermark{#1}}

\gdef\grothFolder{?}\gdef\grothBatch{}\gdef\grothFirst{?}\gdef\grothLast{?}%
\gdef\grothFoldertitle{}\gdef\grothShelfmark{}\gdef\grothDating{}\gdef\grothWatermark{}

% --- The résumé -------------------------------------------------------------
\newenvironment{resume}
  {\small\setlength{\parskip}{0.45em}}
  {\par\medskip\hrule\medskip}

% --- Keywords ---------------------------------------------------------------
% In English, closing the modernised reading's résumé: the modern vocabulary
% under which to search for what the leaves construct. The manifest extracts
% them to tag the volume — this line is the single source of the tags.
\newcommand{\keywords}[1]{\par\smallskip\noindent
  {\footnotesize\textsc{Keywords} --- \textit{#1}}\par}

% --- The critical apparatus -------------------------------------------------

% The Gallica view number, in the margin. This is the anchor that lets a
% disagreement over a formula be settled by looking at *one* image rather than
% a volume of 750. The site uses it to turn the facsimile as the transcription
% is scrolled. A view is one scanned image: usually one side of a leaf.
\newcommand{\page}[1]{%
  \marginnote{\footnotesize\color{gray!70}\textsf{v.\,#1}}%
  \label{page:#1}\ignorespaces}

% The BnF's pencilled foliation, where the leaf carries it and a pass can read
% it: « 198r ». This is what the literature cites — « ff. 198r–208v » — and
% the one line that turns a citation into a view. Recorded, never inferred:
% a folio a pass cannot read is not written.
\newcommand{\folio}[1]{%
  \marginnote{\footnotesize\color{gray!50}\textsf{f.\,#1}}\ignorespaces}

% The modernised reading's coarser counterpart to \page{}: which views a
% section is drawn from.
\newcommand{\pagerange}[2]{%
  \marginnote{\footnotesize\color{gray!70}\textsf{v.\,#1--#2}}%
  \ignorespaces}

% Illegible: never guessed.
\newcommand{\ill}{\textcolor{red!60!black}{[\,\ldots\,]}}

% A doubtful reading: the word is given, and flagged as uncertain.
\newcommand{\uncertain}[1]{\uline{#1}}

% An editorial addition — a missing letter, an implied word.
\newcommand{\add}[1]{\textcolor{blue!50!black}{[#1]}}

% Struck out by the writer. What was crossed out is part of the document.
\newcommand{\struck}[1]{\sout{#1}}

% The transcriber's note, on the state of the page or the mathematics.
\newcommand{\note}[1]{\par\smallskip{\small\color{gray!60!black}\textit{[#1]}}\par\smallskip}

% A marginal note of hers, distinct from the body of the text.
\newcommand{\marginal}[1]{\par\smallskip\noindent\hspace{1em}%
  {\small\color{gray!60!black}#1}\par\smallskip}

% --- Title ------------------------------------------------------------------
% The title block is French, like the documents themselves.

\AddToHook{shipout/background}{%
  \ifx\grothWatermark\empty\else
    \begin{tikzpicture}[remember picture, overlay]
      \node[rotate=52, scale=3.4, align=center, text=gray!18,
            font=\sffamily\bfseries]
        at (current page.center) {\grothWatermark};
    \end{tikzpicture}%
  \fi}

\AtBeginDocument{%
  \begin{center}
    {\large\bfseries Manuscrits de mathématiciens (Gallica) ---
      \ifx\grothShelfmark\empty\grothFolder\else\grothShelfmark\fi}\\[2pt]
    {\itshape\grothFoldertitle}\\[4pt]
    {\small vues \grothFirst--\grothLast
      \ifx\grothBatch\empty\else\ (lot \grothBatch)\fi}\\[2pt]
    \ifx\grothDating\empty\else
      {\small\color{gray!60!black}Datation du catalogue : \grothDating}\\[2pt]
    \fi
    \ifx\grothWatermark\empty\else
      {\small\bfseries\color{red!45!gray}\grothWatermark}\\[2pt]
    \fi
    {\footnotesize\color{gray!60!black}Transcription établie d'après les images IIIF de Gallica.
     Source gallica.bnf.fr / Bibliothèque nationale de France.\\
     Les lectures incertaines sont soulignées, les illisibles notées [\,\ldots\,],
     les ajouts entre crochets ; \textsf{v.} numérote les vues de Gallica,
     \textsf{f.} la foliotation au crayon de la BnF.}
  \end{center}
  \vspace{1em}}

\endinput


\folder{fr-9115}
\shelfmark{Français 9115}
\batch{2}
\pages{21}{40}
\dating{1801-1900}
\watermark{Lecture automatique\\non vérifiée}
\foldertitle{Papiers de Sophie GERMAIN. — Recueil de dissertations et problèmes mathématiques et physiques.}

\begin{document}

\note{Les deux séries de numéros à l'encre relevées au lot 1 se poursuivent
sans interruption : vues 21, 23, 25, 27, 29, 31, 33, 35, 37, 39 → 11, 12, 13,
14, 15, 16, 17, 18, 19, 20, accompagnés chacun d'un second nombre biffé → 19,
21, 23, 25, 27, 29, 31, 33, 35, 37 ; vues 22, 24, 26, 28, 30, 32, 34, 36, 38,
40 → 20, 22, 24, 26, 28, 30, 32, 34, 36, 38. Ces numéros ne sont pas au
crayon : aucune foliotation de la Bibliothèque n'a été lue sur ces vues, et
aucun « f. » n'est donc porté dans cette transcription.}

\page{21}
composé, tendent à se replacer dans la situation qu'une cause extérieure leur
a fait perdre.

L'autre suppose une \emph{tendance des} molécules des corps à se repousser
mutuellement. Celle-ci ne permet donc pas d'attribuer de forces \emph{propres}
à chaque point considéré isolement.

La première, générale comme la définition même des forces d'élasticité,
conduit au terme
\[
  -\left( \frac{1}{\rho} + \frac{1}{\rho'}
          - \left( \frac{1}{R} + \frac{1}{R'} \right) \right)
   \delta\left( \frac{1}{\rho} + \frac{1}{\rho'}
          - \left( \frac{1}{R} + \frac{1}{R'} \right) \right)
\]
et ce terme s'applique, comme on devoit s'y attendre, à toutes les figures
naturelles possibles de la surface.

La seconde ne semble devoir s'appliquer qu'aux seuls cas où les quantités
$\frac{1}{\rho} + \frac{1}{\rho'} - \left( \frac{1}{R} + \frac{1}{R'} \right)$
et $\frac{1}{\rho} - \frac{1}{\rho'}$, sont nulles en même tems.

En effet les forces répulsives ne tendent qu'à diminuer la différence
$\frac{1}{\rho} - \frac{1}{\rho'}$ ; l'action des forces répulsives ne peut
donc être épuisée que quand la différence
$\frac{1}{\rho} - \frac{1}{\rho'}$ devient nulle. Les forces d'élasticité ne
tendent qu'à diminuer la différence
$\frac{1}{\rho} + \frac{1}{\rho'} - \left( \frac{1}{R} + \frac{1}{R'} \right)$ ;
l'action des forces d'élasticité ne peut donc être épuisée que quand la
différence
$\frac{1}{\rho} + \frac{1}{\rho'} - \left( \frac{1}{R} + \frac{1}{R'} \right)$
devient nulle, c'est-à-dire quand la surface a repris sa figure naturelle.
S'il arrivoit que l'équation
$\frac{1}{\rho} + \frac{1}{\rho'} - \left( \frac{1}{R} + \frac{1}{R'} \right) = 0$
n'entraînat pas la condition
$\frac{1}{\rho} - \frac{1}{\rho'} = 0$, les forces répulsives continueroient
d'agir lorsque les forces d'élasticité auroient épuisées leur action. Les deux
hypothèses ne pourroient donc plus s'appliquer à la même nature de surface.
Ainsi, celle qui admet les forces répulsives sera bornée au seul cas où, la
surface considerée dans son état naturel, aura, en chaque point, les rayons
principaux de courbure

\page{22}
égaux et dirigés du même côté.

La condition d'une même limite pour les différences
$\frac{1}{\rho} + \frac{1}{\rho'} - \left( \frac{1}{R} + \frac{1}{R'} \right)$
et $\frac{1}{\rho} - \frac{1}{\rho'}$, ne suffit pas pour que les deux
hypothèses soient équivalentes. Si on réfléchit sur leur nature on verra
qu'une différence essentielle les distingue.

L'une veut que la force élastique soit \emph{propre} a chaque point matériel
qui compose la surface. L'autre suppose que chacun des mêmes points ne doit la
force dont il est doué qu'a l'action qu'il exerce sur les points adjacens et à
celle à laquelle il est soumis de la part des mêmes points.

Les deux hypothèses ne peuvent donc être équivalentes que dans les cas où les
tendances répulsives sont combinées de manière que leur effet sur chaque point
matériel puisse être considéré isolément et comme ne dépendant plus que de
\struck{son} déplacement de ce point.

S'il s'agit, par exemple, d'une simple lame droite, le choix
\uncertain{de l'hypothèse} sera parfaitement indifférent.
\note{« de l'hypothèse » est ajouté dans l'interligne, au-dessus de « le choix
sera »}
Pour nous en convaincre, imaginons trois points matériels $A$, $B$, $C$
disposés en ligne droite et supposons que chacun de ces points exerce sur les
points adjacens une égale action. La force avec laquelle le point $B$ sera
repoussé par le point $A$ sera détruite par la force avec laquelle le même
point $B$ sera repoussé par le point $C$. Le point $B$ ne sera donc plus animé
que par la somme des forces avec lesquelles il repoussera les points adjacens
$A$ et $C$. Le point $A$ sera animé par la force avec laquelle il repoussera le

\page{23}
des trois points se trouvera soumis a la même action que s'ils étoient animés
de forces \emph{propres} uniquement dépendantes de leur déplacement et, quelle
que soit l'étendue de la lame élastique elle sera en équilibre lorsque tous
ses points matériels seront placés a distances égales les uns des autres,
c'est-à-dire lorsqu'elle aura repris sa figure naturelle.

Aussi les différences
$\frac{1}{\rho} + \frac{1}{\rho'} - \left( \frac{1}{R} + \frac{1}{R'} \right)$
et $\frac{1}{\rho} - \frac{1}{\rho'}$, ont-elles, alors, non seulement la même
limite, mais encore la même valeur savoir $\frac{1}{\rho}$.

A l'égard des surfaces naturellement planes le choix de l'hypothèse cesse
d'être indifferent. Cependans aucun des points intérieurs d'une telle surface
se trouve animé par des forces qui ont des sphères d'activité égales entre
elles, ces forces tendent donc a produire le même effet que si elles étoient
réunis a leurs centres respectifs, ou, ce qui revient au même, que si elles
étoient \emph{propres} a chacun des points intérieurs de la surface.
\note{« Cependans aucun… se trouve animé » : la négation attendue (« n'est
animé ») manque sur le feuillet ; la phrase est transcrite telle quelle}

L'ensemble de ces points doit se trouver en équilibre par rapport aux forces
répulsives lorsqu'ils seront tous placés à distance égales les uns des
autres ; c'est-à-dire lorsque la surface aura repris sa figure naturelle.

Suivant la remarque de l'habile géomètre dont j'ai plusieurs fois cité le
mémoire, lorsqu'on a seulement égard aux termes qui après les intégrations

\page{24}
par parties, restent sous le double signe on a identiquement
\[
  \delta \iint \frac{k\, dx\, dy}{\rho\rho'}
  = \iint \delta\!\left( \frac{k}{\rho\rho'} \right) dx\, dy = 0 .
\]
Puisque c'est par le terme $\iint \delta\!\left( \frac{4k}{\rho\rho'} \right) dx\, dy$
que les deux hypothèses diffèrent, il s'en suit que si le choix de
l'hypothèse est indifferent a l'égard des mêmes points, le résultat de l'une
et l'autre hypothèses devienne pour eux identique.
\note{la valeur du coefficient est lue « $4k$ » à la seconde intégrale et
« $k$ » à la première ; les deux lectures sont celles du feuillet}

Il n'en est pas de même pour les points d'extrémités. L'équation à laquelle
ils doivent satisfaire depend du choix de l'hypothèse et, en général, les
conditions auxquelles on peut concevoir que ces points sont assujettis ne
peuvent s'appliquer à la fois aux équations données par l'une et l'autre
hypothèses.

Il est facile d'en voir la raison. Supposons, par exemple, que la plaque
élastique soit terminée par une ligne droite ; pour que les points matériels
situés sur cette ligne soient, par rapport aux forces répulsives, dans le même
cas que les points intérieurs, il faut que la ligne d'extrémité se meuve
parallèlement à elle même, afin que les actions hémisphériques exercés par
chacun des points qui la composent

\page{25}
sur les points adjacens, ajoutés aux actions également hémisphériques que les
points adjacens exercent sur les mêmes points d'extrémités, donne
\uncertain{pour chacun} une action sphérique qui, comme pour les points
intérieurs, puisse être censée réunie à leurs centres et devienne ainsi
équivalente aux forces \emph{propres} que l'autre hypothèse attribue à tous
les points matériels de la surface.
\note{« pour chacun » est ajouté dans l'interligne, au-dessus de « donne une
action »}

N° 7 L'auteur du \emph{mémoire sur les surfaces élastiques} a borné ses
recherches au cas des surfaces naturellement plane. Cet habile géomètre n'a
pas dissimulé les difficultés que l'on rencontreroit dans l'application de son
hypothèse si on vouloit déterminer les conditions auxquelles les points
d'extrémités peuvent être assujettis ; car on lit le passage suivant p. 202
N° 14 de son memoire :

\begin{quote}
« Si le point \emph{m} est, au contraire, situé sur le contour de la surface,
ou qu'il en soit très peu éloigné, la sphère d'activité de la force répulsive,
autour de ce point, ne sera plus complète, c'est-à-dire, qu'une portion de son
étendue ne renfermera pas de points de la surface. Alors les intégrales
relatives à $u$ et $\varphi$ devront être prises dans d'autres limites ; mais
pour trouver l'équation différentielle de la surface élastique en équilibre,
il suffit de considérer les points intérieurs, situés à une distance
quelconque
\end{quote}

\page{26}
\begin{quote}
de son contour ; et l'on n'a besoin d'examiner ce qui arrive aux points
extrêmes, que pour déterminer les forces particulières que l'on doit appliquer
aux limites de la surface pour la tenir en équilibre ; \emph{détermination très
délicate}, sur laquelle je me propose de revenir par la suite mais dont il ne
sera pas question dans ce mémoire. »
\end{quote}

Si l'auteur se fut occupé de la détermination dont il parle ici, j'aurois été
à même de comparer l'une et l'autre théories sous des rapports où elles ne
peuvent manquer de différer essentiellement.

Réduite à mes seules lumières, j'avouerai que l'hypothèse des forces
répulsives me semble devoir mener, sauf le cas dont j'ai parlé N° 6, à
considérer la surface plane comme fesant partie d'une sphère dont le rayon
seroit infini, c'est-à-dire comme infinie elle même.

Il me semble en effet que ce n'est que quand les points doués de forces
répulsives appartiennent à une surface sphérique, que l'on peut toujours
regarder l'effort de chacun de ces points pour reprendre sa situation
naturelle comme égale a l'effort des mêmes points pour se mettre dans la
situation où les forces répulsives se font équilibre ; condition qui me paroît
nécessaire pour que les deux hypothèses soient, généralement parlant,
équivalentes.

\page{27}
Je ne me dissimule pas que le point de vue sous lequel j'ai envisagé la
question est entièrement neuf ; et c'est pour moi un grand sujet de défiance.

Dans le cas linéaire, les géomètres ont eu recours à la considération de
l'angle de contingence, et ils ont présentés de grandes difficultés à
exprimer, pour la surface, un nombre infini de flexions qui s'exécuteroient à
la fois dans tous les sens. On peut regarder l'hypothèse qui admet la somme
des raisons inverses des deux rayons de principales courbures, comme
introduisant les angles de contingences des deux courbures respectives ; mais
je vais plus loin et je ne regarde les variations de courbures que comme le
resultat de la variation de position de chacun des points matériels dont se
compose la surface. Je crois me conformer en cela a la définition des forces
d'élasticité ; car le corps élastique ne peut tendre à reprendre sa forme et
son étendue naturelles qu'autant que les points materiels qui le composent
tendront eux mêmes à reprendre leur position naturelle. Cette idée est
fondamentale ; elle réduit la question à ses élemens les plus simples, je
supplie les géomètres de lui accorder quelqu'attention.

\page{28}
N° 8 Suivant la remarque du savant auteur du \emph{Mémoire sur les surfaces
élastiques} (voyez N° 23 p. \uncertain{221}) la plaque élastique en état
d'équilibre est parmi toutes les surfaces de même étendue, celle dans laquelle
l'intégrale
$\iint \left( \frac{1}{\rho} + \frac{1}{\rho'} \right)^{2} k\, dx\, dy$
est un \emph{maximum} ou un \emph{minimum}.

Si j'ai reussi à prouver que le terme dû aux forces d'élasticité est pour
chaque point de la surface
$-\left( \frac{1}{\rho} + \frac{1}{\rho'} - \left( \frac{1}{R} + \frac{1}{R'} \right) \right)
 \delta\left\{ \frac{1}{\rho} + \frac{1}{\rho'} - \left( \frac{1}{R} + \frac{1}{R'} \right) \right\}$ ;
la même remarque s'étendra a la surface élastique de figure quelconque, et on
pourra affirmer que cette surface est, parmi toutes les surfaces de même
étendue, celle pour laquelle l'intégrale
\[
  \iint \left\{ \frac{1}{\rho} + \frac{1}{\rho'}
        - \left( \frac{1}{R} + \frac{1}{R'} \right) \right\}^{2} k\, dx\, dy
\]
est un \emph{maximum} ou un \emph{minimum}.

J'espere être, par la suite, en état de démontrer que ce théorème conduit en
effet a l'équation de la surface élastique de figure quelconque

Dans les §§ suivans je ne m'occuperai plus que du seul cas des plaques
vibrantes

\page{29}
\section{Conclusions analytiques déduites de l'une et l'autre hypothèses}

N° 9 Mon hypothèse conduit au terme
$-\left\{ \frac{1}{\rho} + \frac{1}{\rho'}
  - \left( \frac{1}{R} + \frac{1}{R'} \right) \right\}
 \delta\left\{ \frac{1}{\rho} + \frac{1}{\rho'}
  + \left( \frac{1}{R} + \frac{1}{R'} \right) \right\}$
en negligeant les quantités du second ordre, ce terme se réduit ici, à cause
de $\frac{1}{R} + \frac{1}{R'} = 0$, à
$-\left( \frac{d^{2}z}{dx^{2}} + \frac{d^{2}z}{dy^{2}} \right)
 \delta\left( \frac{d^{2}z}{dx^{2}} + \frac{d^{2}z}{dy^{2}} \right)$.
\note{à la seconde accolade la vue porte un signe $+$ devant
$\left( \frac{1}{R} + \frac{1}{R'} \right)$, là où la vue 28 écrit deux fois
$-$ dans le même terme}

Étendu à la surface entière il donne l'intégrale double
\[
  -\iint \left\{ \left( \frac{d^{2}z}{dx^{2}} + \frac{d^{2}z}{dy^{2}} \right)
    \delta\left( \frac{d^{2}z}{dx^{2}} + \frac{d^{2}z}{dy^{2}} \right)
    \right\} dx\, dy .
\]

On a en intégrant par parties
\begin{align*}
  &\iint \left\{ \left( \frac{d^{2}z}{dx^{2}} + \frac{d^{2}z}{dy^{2}} \right)
    \delta\left( \frac{d^{2}z}{dx^{2}} + \frac{d^{2}z}{dy^{2}} \right)
    \right\} dx\, dy \\
  &\quad = \int \left( \frac{d^{2}z}{dx^{2}} + \frac{d^{2}z}{dy^{2}} \right)
      \delta\frac{dz}{dx} . dy
   - \int \frac{d\left\{ \frac{d^{2}z}{dx^{2}} + \frac{d^{2}z}{dy^{2}} \right\}}{dx}
      \delta z\, dy \\
  &\qquad + \int \left( \frac{d^{2}z}{dx^{2}} + \frac{d^{2}z}{dy^{2}} \right)
      \delta\frac{dz}{dy} . dx
   - \int \frac{d\left\{ \frac{d^{2}z}{dx^{2}} + \frac{d^{2}z}{dy^{2}} \right\}}{dy}
      \delta z\, dx \\
  &\qquad + \iint \left\{ \frac{d^{4}z}{dx^{4}}
      + 2\frac{d^{4}z}{dx^{2}dy^{2}} + \frac{d^{4}z}{dy^{4}} \right\}
      \delta z\, dx\, dy .
\end{align*}

N° 10 L'hypothèse des forces répulsives conduit au terme
$-\left\{ \frac{1}{\rho} - \frac{1}{\rho'} \right\}
 \delta\left\{ \frac{1}{\rho} - \frac{1}{\rho'} \right\}$
en négligeant les quantités du second ordre, ce terme se réduit ici à
\begin{gather*}
  -\left\{ \sqrt{ \left( \frac{d^{2}z}{dx^{2}} + \frac{d^{2}z}{dy^{2}} \right)^{2}
     - 4\left\{ \frac{d^{2}z}{dx^{2}} . \frac{d^{2}z}{dy^{2}}
     - \left( \frac{d^{2}z}{dx\,dy} \right)^{2} \right\} } \right\} \\
  \delta\left\{ \sqrt{ \left( \frac{d^{2}z}{dx^{2}} + \frac{d^{2}z}{dy^{2}} \right)^{2}
     - 4\left\{ \frac{d^{2}z}{dx^{2}} . \frac{d^{2}z}{dy^{2}}
     - \left( \frac{d^{2}z}{dx\,dy} \right)^{2} \right\} } \right\} \\
  = -\left( \frac{d^{2}z}{dx^{2}} + \frac{d^{2}z}{dy^{2}} \right)
     \delta\left( \frac{d^{2}z}{dx^{2}} + \frac{d^{2}z}{dy^{2}} \right) \\
   - 2\left\{ \frac{d^{2}z}{dx^{2}} \delta\frac{d^{2}z}{dy^{2}}
     + \frac{d^{2}z}{dy^{2}} \delta\frac{d^{2}z}{dx^{2}}
     - 2\frac{d^{2}z}{dx\,dy} \delta\frac{d^{2}z}{dx\,dy} \right\} .
\end{gather*}
\note{le coefficient de l'accolade est lu $-2$ ; la différentiation du radical
donnerait $+2$, mais le feuillet porte le signe négatif, ici comme à la
vue 30}

Étendu à la surface entière il donne l'intégrale double
\[
  -\iint \left\{ \left( \frac{d^{2}z}{dx^{2}} + \frac{d^{2}z}{dy^{2}} \right)
     \delta\left( \frac{d^{2}z}{dx^{2}} + \frac{d^{2}z}{dy^{2}} \right)
   - 2\left\{ \frac{d^{2}z}{dx^{2}} \delta\frac{d^{2}z}{dy^{2}}
     + \frac{d^{2}z}{dy^{2}} \delta\frac{d^{2}z}{dx^{2}}
     - 2\frac{d^{2}z}{dx\,dy} \delta\frac{d^{2}z}{dx\,dy} \right\} \right\}
   dx\, dy
\]
en intégrant par parties on trouve :

\page{30}
\begin{align*}
  &\iint \left\{ \frac{d^{2}z}{dx^{2}} \delta\frac{d^{2}z}{dy^{2}}
    + \frac{d^{2}z}{dy^{2}} \delta\frac{d^{2}z}{dx^{2}}
    - 2\frac{d^{2}z}{dx\,dy} \delta\frac{d^{2}z}{dx\,dy} \right\} dx\, dy \\
  &\quad = \int \frac{d^{2}z}{dx^{2}} \delta\frac{dz}{dy} . dx
   - \int \frac{d^{3}z}{dx^{2}dy} \delta z\, dx
   + \iint \frac{d^{4}z}{dx^{2}dy^{2}} \delta z . dx\, dy \\
  &\qquad + \int \frac{d^{2}z}{dy^{2}} \delta\frac{dz}{dx} . dy
   - \int \frac{d^{3}z}{dy^{2}dx} \delta z . dy
   + \iint \frac{d^{4}z}{dx^{2}dy^{2}} \delta z . dx\, dy \\
  &\qquad - 2\left\{ -\int \frac{d^{2}z}{dx\,dy} \delta\frac{dz}{dx} . dx
   - \int \frac{d^{3}z}{dy^{2}dx} . \int \delta\ill{}\, dy
   + \iint \frac{d^{4}z}{dx^{2}dy^{2}} \delta z\, dx\, dy \right\}
   + \frac{d^{2}z}{dx\,dy} \delta z .
\end{align*}
\note{dans la dernière accolade un terme est biffé puis récrit par-dessus, et
ni la couche biffée ni la correction ne se laissent lire : la vue porte
« $\int \delta$ », un mot rayé, puis « $dy$ », avec « $dy$ » repris au-dessous.
Le terme « $+\,\frac{d^{2}z}{dx\,dy}\delta z$ » est ajouté sur une ligne
propre, sous l'accolade}

On a par conséquent en ayant égard à la valeur de
$\iint \left\{ \left( \frac{d^{2}z}{dx^{2}} + \frac{d^{2}z}{dy^{2}} \right)
\delta\left( \frac{d^{2}z}{dx^{2}} + \frac{d^{2}z}{dy^{2}} \right) \right\}
dx\, dy$
donnée N° précédent et en effaçant ce qui se détruit :
\begin{align*}
  &\iint \left\{ \left( \frac{d^{2}z}{dx^{2}} + \frac{d^{2}z}{dy^{2}} \right)
     \delta\left( \frac{d^{2}z}{dx^{2}} + \frac{d^{2}z}{dy^{2}} \right)
   - 2\left\{ \frac{d^{2}z}{dx^{2}} \delta\frac{d^{2}z}{dy^{2}}
     + \frac{d^{2}z}{dy^{2}} \delta\frac{d^{2}z}{dx^{2}}
     - 2\frac{d^{2}z}{dx\,dy} \delta\frac{d^{2}z}{dx\,dy} \right\} \right\}
     dx\, dy \\
  &\quad = \int \left( \frac{d^{2}z}{dx^{2}} + \frac{d^{2}z}{dy^{2}} \right)
      \delta\frac{dz}{dx} . dy
   - \int \frac{d\left\{ \frac{d^{2}z}{dx^{2}} + \frac{d^{2}z}{dy^{2}} \right\}}{dx}
      \delta z\, dy \\
  &\qquad + \int \left( \frac{d^{2}z}{dx^{2}} + \frac{d^{2}z}{dy^{2}} \right)
      \delta\frac{dz}{dy} . dx
   - \int \frac{d\left\{ \frac{d^{2}z}{dx^{2}} + \frac{d^{2}z}{dy^{2}} \right\}}{dy}
      \delta z\, dx \\
  &\qquad - 2\left\{ \int \frac{d^{2}z}{dy^{2}} \delta\frac{dz}{dx} . dy
      \ill{} \int \frac{d^{3}z}{dy^{2}dx} \delta z\, dy \right\} \\
  &\qquad - 2\left\{ \int \frac{d^{2}z}{dx^{2}} \delta\frac{dz}{dy} . dx
      \struck{- \int \frac{d^{3}z}{dx^{2}dy} \delta\frac{dz}{dx} . dx}
      \struck{- \int \frac{d^{3}z}{dx^{2}dy} \delta z\, dx} \right\} \\
  &\qquad + 4\left\{ \frac{d^{2}z}{dx\,dy} \delta z
      - \int \frac{d^{2}z}{dx\,dy} \delta\frac{dz}{dx} dx
      - \int \frac{d^{2}z}{dx\,dy} \delta\frac{dz}{dy} dy \right\} \\
  &\qquad + \iint \left\{ \frac{d^{4}z}{dx^{4}}
      + 2\frac{d^{4}z}{dx^{2}dy^{2}} + \frac{d^{4}z}{dy^{4}} \right\}
      \delta z\, dx\, dy .
\end{align*}
\note{ces deux dernières accolades sont les plus retouchées du lot : le signe
qui joint les deux intégrales de la première est couvert d'un pâté d'encre et
n'a pas été lu ; dans la seconde, deux groupes sont rayés, et la lecture de la
couche annulée reste douteuse ; le $2$ de « $-2$ » est repassé à l'encre}

N° 11 Conformement à ce qui a été dit N° 6, les termes qui restent sous le
double signe, c'est-à-dire ceux qui se rapportent aux points intérieurs sont
indépendans du choix de l'une ou l'autre hypothèses.

Si on multiplie la somme de ces termes par le coefficient constant $n^{2}$
proportionnel à l'épaisseur et a l'élasticité naturelle de la surface, et que
l'on ajoute le terme

\page{31}
$\frac{d^{2}z}{dt^{2}}$ dû aux forces accélératrices, on aura, en ayant égard
au signe
\[
  \frac{d^{2}z}{dt^{2}} + n^{2}\left( \frac{d^{4}z}{dx^{4}}
  + 2\frac{d^{4}z}{dx^{2}dy^{2}} + \frac{d^{4}z}{dy^{4}} \right) = 0
\]
pour l'équation de la plaque élastique vibrante.

C'est en effet cette équation qui avoit été calculée d'après mon hypothèse et
qui se trouve dans la pièce que j'ai envoyée pour le concours de janvier 1814.

A l'égard des termes qui se rapportent aux extremités calculés suivant l'une
et l'autre hypothèses, ils ne peuvent, en général, être regardés comme
équivalens. Cependans si on vouloit que la surface plane fut infinie, on
auroit $x$ et $y$ pareillement infinis et par conséquent $dx = 0$, $dy = 0$ ;
les termes qui se rapportent aux extremités disparoîtroient d'eux mêmes et,
comme je l'ait dit N° 7, le choix de l'hypothèse seroit parfaitement
indifferent.

Dans tout autre cas il me reste a examiner commens la différence entre les
mêmes termes s'accomode avec l'experience.

\section{Comparaison entre les resultats du calcul et les phénomènes observés}

N° 12 Mon but actuel n'etant pas de traiter la théorie des surfaces
élastiques, mais seulement de mettre le lecteur à même d'apréciér les
differentes hypothèses employés à la recherche de l'équation

\page{32}
de ce genre de surfaces ; je n'indiquerai parmi les phénomènes qui ont été
expliqués dans les mémoires que j'ai présentés a l'accademie, que ceux qui ont
un rapport direct avec le choix de l'hypothèse.

J'éviterai le plus possible de citer les faits particuliers, et d'entrer dans
le détail de leur explication.

L'approbation de l'institut est une garantie suffisante de l'accord entre
l'experience et la théorie déduite de mon hypothèse. Je partirai donc de cet
accord comme d'un point convenu ; et, en bornant les exemples à un petit
nombre, je tacherai de faire apréciér la différence des resultats obtenus du
calcul, suivant que l'on employe l'une ou l'autre hypothèses proposées.

N° 13 En suivant mon hypothèse nous avons trouvé, N° 9, que les termes
relatifs aux limites sont :
\begin{align*}
  &\int \left( \frac{d^{2}z}{dx^{2}} + \frac{d^{2}z}{dy^{2}} \right)
      \delta\frac{dz}{dx} . dy
   - \int \frac{d\left\{ \frac{d^{2}z}{dx^{2}} + \frac{d^{2}z}{dy^{2}} \right\}}{dx}
      \delta z\, dy \\
  &\quad + \int \left( \frac{d^{2}z}{dx^{2}} + \frac{d^{2}z}{dy^{2}} \right)
      \delta\frac{dz}{dy} . dx
   - \int \frac{d\left\{ \frac{d^{2}z}{dx^{2}} + \frac{d^{2}z}{dy^{2}} \right\}}{dy}
      \delta z\, dx
\end{align*}

L'hypothèse des forces répulsives nous a donné, N° 10, aussi pour les limites
les termes suivans :
\begin{align*}
  &\int \left( \frac{d^{2}z}{dx^{2}} + \frac{d^{2}z}{dy^{2}} \right)
      \delta\frac{dz}{dx} . dy
   - \int \frac{d\left\{ \frac{d^{2}z}{dx^{2}} + \frac{d^{2}z}{dy^{2}} \right\}}{dx}
      \delta z\, dy \\
  &\quad + \int \left( \frac{d^{2}z}{dx^{2}} + \frac{d^{2}z}{dy^{2}} \right)
      \delta\frac{dz}{dy} . dx
   - \int \frac{d\left\{ \frac{d^{2}z}{dx^{2}} + \frac{d^{2}z}{dy^{2}} \right\}}{dy}
      \delta z\, dx
\end{align*}

\page{33}
\begin{align*}
  &-2\left\{ \int \frac{d^{2}z}{dy^{2}} \delta\frac{dz}{dx} . dy
     + \int \frac{d^{3}z}{dy^{2}dx} \delta z\, dy \right\} \\
  &\quad - 2\left\{ \int \frac{d^{2}z}{dx^{2}} \delta\frac{dz}{dy} . dx
     - 2\int \frac{d^{2}z}{dx\,dy} \delta\frac{dz}{dx} . dx
     - \int \frac{d^{3}z}{dx^{2}dy} \delta z\, dx \right\}
\end{align*}
ou, en effaçant ce qui se détruit et mettant au lieu de
$\int \frac{d^{2}z}{dx\,dy} \delta\frac{dz}{dx} dx$ sa valeur
$\frac{d^{2}z}{dx\,dy} \delta z - \int \frac{d^{3}z}{dx^{2}dy} \delta z\, dx$ :
\begin{align*}
  &2\frac{d^{2}z}{dx\,dy} \delta z
   + \int \left( \frac{d^{2}z}{dx^{2}} - \frac{d^{2}z}{dy^{2}} \right)
      \delta\frac{dz}{dx} . dy
   - \int \frac{d\left\{ \frac{d^{2}z}{dx^{2}} \ill{} \frac{d^{2}z}{dy^{2}} \right\}}{dx}
      \delta z\, dy \\
  &\quad - 4\int \frac{d^{2}z}{dx\,dy} \delta\frac{dz}{dx}\, dx
   + \int \left( \frac{d^{2}z}{dy^{2}} - \frac{d^{2}z}{dx^{2}} \right)
      \delta\frac{dz}{dy} . dx
   - \int \frac{d\left\{ \frac{d^{2}z}{dy^{2}} \ill{} \frac{d^{2}z}{dx^{2}} \right\}}{dy}
      \delta z\, dx \\
  &\quad - 4\int \frac{d^{2}z}{dx\,dy} \delta\frac{dz}{dy}\, dy
\end{align*}
\note{aux deux accolades différentiées le signe qui joint les deux dérivées
secondes est couvert d'un pâté d'encre et n'a pas été lu ; les parenthèses qui
précèdent, où le même couple paraît, portent un signe $-$}

Supposons qu'il s'agisse d'une plaque parallélogramme et que $x$, par exemple,
soit constant pour tous les points qui appartiennent à la ligne d'extrémité ;
suivant mon hypothèse l'équation
\[
  \left( \frac{d^{2}z}{dx^{2}} + \frac{d^{2}z}{dy^{2}} \right) \delta\frac{dz}{dx}
  - \frac{d\left\{ \frac{d^{2}z}{dx^{2}} + \frac{d^{2}z}{dy^{2}} \right\}}{dx}
    \delta z = 0 .
\]
aura lieu pour tous ces points. On pourra y satisfaire en supposant à la fois
ou, $\frac{d^{2}z}{dx^{2}} + \frac{d^{2}z}{dy^{2}} = 0$ et $\delta z = 0$ ;
ou,
$\frac{d\left\{ \frac{d^{2}z}{dx^{2}} + \frac{d^{2}z}{dy^{2}} \right\}}{dx} = 0$
et $\delta\frac{dz}{dx} = 0$.

Dans le mémoire envoyé pour le concours de janvier 1814 j'ai expliqué un grand
nombre de figures et de sons correspondans, à l'aide de l'intégrale
\struck{$z = \ldots\, \sin\frac{my\pi}{A} \sin$}
\[
  z = \ldots\, \sin\frac{my\pi}{A} \sin . \frac{nx\pi}{A} .
\]
$m$ et $n$ étant des nombres entiers il est évident qu'aux extrémités
c'est-à-dire lorsque $x = A$ ou $x = 0$ on a $\sin\frac{nx\pi}{A} = 0$ et par
conséquent aussi $\frac{d^{2}z}{dx^{2}} + \frac{d^{2}z}{dy^{2}} = 0$ et
$z = 0$.

\page{34}
Lorsque $x = \frac{A}{2}$ on a $\cos\left( \frac{n x \pi}{A} \right)
= \cos . n . \frac{\pi}{2} = 0$ et par conséquent aussi
$\frac{dz}{dx} = 0$ et
$\frac{d\left\{ \frac{d^{2}z}{dx^{2}} + \frac{d^{2}z}{dy^{2}} \right\}}{dx} = 0$.

La surface est partagée par des lignes d'extrèmités qui se croissent à angles
droits. Celles qui satisfont aux conditions $z = 0$ et
$\frac{d^{2}z}{dx^{2}} + \frac{d^{2}z}{dy^{2}} = 0$ sont les lignes nodales de
la surface. Celles qui satisfont aux conditions $\frac{dz}{dx} = 0$ et
$\frac{d\left\{ \frac{d^{2}z}{dx^{2}} + \frac{d^{2}z}{dy^{2}} \right\}}{dx} = 0$
se meuvent sans cesser d'être parallèles à leur position naturelle ; c'est ce
qu'il est aisé de sentir en remarquant que, la position initiale de la plaque
étant supposée horisontale, $\frac{dz}{dx}$ exprime la tangente de l'angle que
la ligne d'extrémité fait avec l'horison. Ces dernières lignes sont appellées
libres par analogie aux points d'extrémités libres dont Euler s'est occupé
dans son mémoire sur la lame.

Les lignes dont je viens de parler peuvent être situées dans l'intérieure de
la plaque : leur nombre dépend des valeurs de \emph{m} et de \emph{n}. Elles
jouissent de cette propriété que si la plaque est coupée suivant ces lignes,
les parties separées donnent le même son que la plaque entière et présentent
des figures nodales dont la réunion reproduit la figure nodale observée sur la
plaque avant sa division.

\page{35}
N° 13 En conservant les suppositions précédentes, si on adopte l'hypothèse des
forces répulsives l'équation
\note{le numéro 13 est employé deux fois : la vue 32 ouvre déjà un « N° 13 »,
consacré à l'hypothèse de l'auteur ; celui-ci, consacré à l'hypothèse des
forces répulsives, porte le même numéro sur le feuillet}
\begin{gather*}
  2\frac{d^{2}z}{dx\,dy} \delta z
  + \int \left( \frac{d^{2}z}{dx^{2}} - \frac{d^{2}z}{dy^{2}} \right)
     \delta\frac{dz}{dx}\, dy
  - \int \frac{d\left\{ \frac{d^{2}z}{dx^{2}} \ill{} \frac{d^{2}z}{dy^{2}} \right\}}{dx}
     \delta z . dy = 0 \\
  - 4\int \frac{d^{2}z}{dx\,dy} \delta\frac{dz}{dy}\, dy
\end{gather*}
\note{le « $=0$ » est porté à la fin de la première ligne, le dernier terme
étant renvoyé au-dessous ; le signe entre les deux dérivées secondes de
l'accolade différentiée est couvert d'un pâté d'encre}
devra avoir lieu pour tous les points qui appartiennent à la ligne
d'extrémité ; \struck{et en effet} les conditions $\delta z = 0$ et
$\frac{d^{2}z}{dx^{2}} + \frac{d^{2}z}{dy^{2}} = 0$ se réduisant ici à
$\sin\left( \frac{n x \pi}{A} \right) = 0$ entraîneront la condition
$\frac{d^{2}z}{dx^{2}} - \frac{d^{2}z}{dy^{2}} = 0$. \struck{De même aussi}
mais elles ne feront pas disparoître le terme
$4\int \frac{d^{2}z}{dx\,dy} \delta\frac{dz}{dy}\, dy$.
\note{« mais elles ne feront pas disparoître le terme… » est écrit dans
l'interligne et au-dessus de la ligne suivante, en remplacement de « De même
aussi », rayé}

Les conditions $\delta\frac{dz}{dx} = 0$,
$\frac{d\left\{ \frac{d^{2}z}{dx^{2}} + \frac{d^{2}z}{dy^{2}} \right\}}{dx} = 0$
se réduisant a $\cos\left( \frac{n x \pi}{A} \right) = 0$ entraîneront les
conditions
$\frac{d\left\{ \frac{d^{2}z}{dx^{2}} + 3\frac{d^{2}z}{dy^{2}} \right\}}{dx} = 0$
et $\frac{d^{2}z}{dx\,dy} = 0$.

Ainsi, dans les cas qui se rapportent à l'intégrale
$z = \ldots\, \sin\left( \frac{m y \pi}{A} \right)
\sin\left( \frac{n x \pi}{A} \right)$,
\struck{cas dans lesquels la ligne de repos ne satisfait plus aux conditions}
des extrémités ; mais pour celle qui se meut parallèlement à sa situation
initiale, les équations d'extrémités données par l'une et l'autre hypothèses
sont satisfaites par les mêmes conditions ; et, comme je l'ai annoncé N° 6,
\struck{le choix de l'hypothèse est indifférent parce} qu'elles deviennent
alors entièrement équivalentes \struck{et ne sont plus que deux manières
différentes d'envisager un seul et même objet}.
\note{tout ce passage est remanié en deux couches, avec des mots rayés et des
corrections dans l'interligne ; la lecture de la couche annulée est donnée
pour ce qu'elle vaut, et le texte retenu est celui qui subsiste}

N° 14 L'intégrale linéaire
\[
  z = \ldots\, \left(
   A\left( a\, e^{\frac{\omega x}{A}} + b\, e^{-\frac{\omega x}{A}}
     + c \sin\frac{x\omega}{A} + d \cos\frac{x\omega}{A} \right)
   + B\left( a\, e^{\frac{\omega y}{A}} + b\, e^{-\frac{\omega y}{A}}
     + c \sin\frac{y\omega}{A} + d \cos\frac{y\omega}{A} \right) \right)
\]

\page{36}
m'a servi à expliquer comment, le son restant le même, les figures nodales les
plus bizares et les plus variées remplaçent souvent, sur la plaque quarrée les
simples lignes droites dont la théorie indique la formation.

Le son et le nombre des lignes droites nodales dependent de la valeur de
$\omega$. Si on fait successivement $A = 0$ et $B = 0$, la figure nodale sera
composée des lignes droites tracées parallèlement à l'une ou l'autre des
lignes d'extremités. L'intersection des lignes nodales formées dans l'une et
l'autre suppositions $A = 0$, $B = 0$ détermine sur la surface des points de
repos \emph{invariables}.

A chaque valeur relative de $A$ et de $B$ correspond une nouvelle figure
nodale, les points \emph{invariables} appartiennent à chacune d'elles ; elles
affectent d'ailleurs les formes les plus variées. Tantôt elles sont composées
de simples lignes sinueuses en même nombre que les lignes droites qu'elles
remplacent ; tantôt elles offrent des lignes diagonales des courbes rentrantes
\&c. Quelquefois même la figure nodale qui se forme sur une moitié de la
plaque appartient à des valeurs differentes de $A$ et $B$ que celle qui se
montre sur l'autre moitié. Cette variété si étonnante au premier coup d'œil
reçoit cependant une explication satisfesante.

Les lignes nodales passent évidemment par tous les

\page{37}
points que satisfont à l'equation $z = 0$, ou, ce qui est la même chose à
l'équation :
\[
  A\left\{ a\, e^{\frac{\omega x}{A}} + b\, e^{-\frac{\omega x}{A}}
    + c \sin . \frac{\omega x}{A} + d \cos . \frac{\omega x}{A} \right\}
  + B\left\{ a\, e^{\frac{\omega y}{A}} + b\, e^{-\frac{\omega y}{A}}
    + c \sin . \frac{\omega y}{A} + d \cos . \frac{\omega y}{A} \right\} = 0 .
\]
Pour les points \emph{invariables} la condition $z = 0$ est satisfaite en
vertu des deux équations :
\begin{gather*}
  a\, e^{\frac{\omega x}{A}} + b\, e^{-\frac{\omega x}{A}}
   + c \sin . \frac{\omega x}{A} + d \cos . \frac{\omega x}{A} = 0 \\
  a\, e^{\frac{\omega y}{A}} + b\, e^{-\frac{\omega y}{A}}
   + c \sin . \frac{\omega y}{A} + d \cos . \frac{\omega y}{A} = 0
\end{gather*}

Dans tout autre cas les valeurs de
$a\, e^{\frac{\omega x}{A}} + e^{-\frac{\omega x}{A}} + \&c$,
$a\, e^{\frac{\omega y}{A}} + e^{-\frac{\omega y}{A}} + \&c$, et par
conséquent aussi celle de $x$ et $y$, qui satisfont a la condition
$z \ill{} 0$, dépendent nécessairement de la valeur relative des constantes
$A$ et $B$.
\note{le signe entre $z$ et $0$ est couvert d'un pâté d'encre et n'a pas été
lu ; la même condition est écrite « $z = 0$ » quelques lignes plus haut}

Si on adopte mon hypothèse la séparation des variables $x$ et $y$ qui présente
notre \struck{équation} intégrale comme une généralisation de celle qui
convient au cas linéaire n'a rien qui doive étonner ; car quelque soit la
forme a laquelle la surface ait été ployée par l'action d'une cause
extérieure, elle tend a reprendre la forme qu'elle affectoit avant l'action
supposée.

Mais si on veut au contraire que l'elasticité soit due à des forces
répulsives, comment concevoir l'application d'une pareille intégrale ? Les
forces repulsives n'agissent elles pas sur tous les points situés dans leur
sphère d'activité, comment donc ces forces deviendroient-elles ici linéaires
et n'agiroient-elles plus que dans des directions dont les fils croisés
imaginés par Euler a l'occasion des surfaces tendues, donneroient une juste
idée ?

\page{38}
N° 15 J'ajouterai encore quelques observations relatives à un phénomène que
j'ai déja fait connoître (M. pour janvier 1814)

Si on applique un corps mou à un des bords d'une plaque élastique en observant
de distribuer également la masse sur tous les points de la ligne
d'application, et qu'on fasse vibrer ensuite la plaque ainsi chargée ; on
obtiendra la même figure et le même son que si la même plaque étoit riellement
prolongée d'une quantité égale en poids au corps ajouté. Il arrive delà que,
plus on augmente ce poids, plus aussi la ligne voisine de l'extrèmité chargée
s'en trouve rapprochée.

Je suis parvenue, comme je l'ai dit dans le mémoire déja cité, à obtenir un
son juste, d'une plaque assez chargée pour que l'intervalle entre deux lignes
nodales intérieures fut dix fois plus grand que la partie de la même plaque
comprise entre l'extrèmité chargée et la ligne nodale la plus voisine.

La manière de vibrer qui est l'objet de cette experience appartient à
l'intégrale $z = \ldots\, \sin\left( \frac{m y \pi}{A} \right)
\sin\left( \frac{n x \pi}{A} \right)$ ; les lignes nodales sont toutes alors
des lignes d'extrèmités analytiques, entorte, que si la plaque étoit
riellement séparée en autant de parties, et qu'on parvient à leur imprimer le
même ébranlement que lorsqu'elles sont reunies ; le son que chacune d'elles
rendroient

\page{39}
seroit le même qu'avant leur séparation. On a donc le singulier phénomène
d'une partie dix fois moins grande qu'une autre dont le son correspondant est
pourtant le même.

J'ai eu depuis occasion d'observer le même fait sur des surfaces courbes. J'ai
toujours trouvé que le poids de la cire collante que j'appliquois à une des
extrèmités de la surface, fesoit varier la figure et le son, des mêmes
quantités que si cette surface étoit riellement prolongée dans le sens de la
dimension chargée, d'une étendue égale en poids, à la cire ajoutée.

L'influence de la substitution se réduit à rendre le son moins intense et les
lignes nodales plus larges. Ces deux effets doivent être attribués a la
résistance qu'une substance tenace oppose naturellement à la vibration de la
surface. Ils ne peuvent manquer d'avoir lieu à la fois ; car lorsque le son est
très intense les lignes nodales sont fort nettes parceque le mouvement des
parties voisines devant alors sensible.
\note{la phrase reste inachevée sur le feuillet : rien ne suit « devant alors »
avant le mot « sensible », porté seul à la ligne}

A l'aide de mon hypothèse, le phénomène dont je viens de parler recoit une
explication fort naturelle. Une surface d'une grandeur déterminée, lorsqu'elle
se prête à une vibration donnée, est susceptible de rendre un certain son et
de présenter une figure nodale correspondante. Aucun des points matériels qui
la composent est alors doué d'une force propre

\page{40}
en vertu de laquelle il tend à reprendre sa position naturelle. Si durant le
mouvement une portion de la surface élastique est remplacée par une matière
inerte, chacun des points élastiques conservés garde la même tendance à
reprendre sa position naturelle, et la surface continue de se mouvoir comme
auparavant. Les mouvemens alternatifs sont à la vérité plus foibles, non
seulemens parceque la matière inerte n'y participe pas, mais encore parceque
sa presence gène la vibration

Si on veut, au contraire, que l'élasticité soit due aux forces répulsives,
comment concevoir qu'il ne s'opere aucun changement essentiel lorsque des
points matériels auxquels de pareilles forces ne sauroient être attribuées
remplacent, par leur poids, les points qui sont doués de ces forces ?

Il s'est élevé plus d'une fois entre les géometres des discussions relatives
aux principes métaphysiques admis dans le calcul. Souvent alors, le lecteur,
qui négligeoit d'examiner les raisons apportés de part et d'autre, avoit à
balancer l'autorité de deux noms célèbres.

Un pareil prestige n'est pas de mon côté, et loin delà, il est contre moi ; je
n'ai donc pas à craindre que, si je me trompe, on admette mon erreur sans
examen. J'ai bien plûtôt a craindre que les géometres justement avares de leur
tems

\note{la vue 40 s'arrête au milieu de la phrase ; le texte se poursuit à la
vue 41, hors de ce lot}

\end{document}
